\section{The full polynomial period deformation and its entire-amplitude forcing}
\label{GraphV3:sec:pde-full-deformation}

Retain the original nonempty complete packet $h(s)$, its degree $n\geq1$,
its primitive $\Phi_h(0)=0$, its arithmetic coefficient space
$E_h=\mathbb{C}[s]/(h)$ in the increasing frame $1,s,\ldots,s^{n-1}$,
and the actual entire amplitude $v_h=2\xi/h$. All zero orders in $h$
are complete, so this last quotient is entire and its Taylor unit
$\upsilon_h=j_h(v_h)$ is invertible in $E_h$. Write
\[
 h(s)=s^n+\sum_{r=0}^{n-1}h_rs^r,\qquad
 f_t(s)=\Phi_h(s)-ts,\qquad
 A_h(t)=A_h+tR,\quad R=e_0e_{n-1}^{\mathsf T}.
 \tag{PDE1}
\]
Here $A_h$ is the original arithmetic multiplication by $s$ on $E_h$;
$A_h(t)$ is its displayed companion deformation on the same coefficient
frame. The parameter $t$ does not replace the original marked arithmetic
quotient by a different quotient. The matrix $R$ includes the case
$n=1$, where it is the identity.

Fix $u\ne0$ on a specified branch of $\log u$. Set $d=n+1$ and retain
the actual rays and their orientations
\[
 \theta_j=\frac{\arg u+\pi+2\pi j}{d},\quad
 \ell_j(r)=re^{i\theta_j}\ (r\geq0),\quad
 \Gamma_j=-\ell_0+\ell_j\quad(1\leq j\leq n).
 \tag{PDE2}
\]
For an admitted amplitude $a$, let $\mathcal I_t(a)$ be the column with
entries $\int_{\Gamma_j}e^{f_t/u}a(s)\,ds$. Every amplitude used below is
a polynomial times a derivative of the fixed $v_h$, or a polynomial.
IPM1--6 proves their convergence on these rays with the original theta
bound; the explicit derivative estimate below also proves the required
first-derivative convergence directly. Compact sets of parameters have
a common dominating function with negative leading term
$-r^d/(d|u|)$ and lower polynomial terms. Thus these integrals and their
displayed parameter derivatives are entire in $t$. The fixed rays have
no $t$-dependent endpoint or orientation.

\subsection{The exact determinant on the whole polynomial family}

For this calculation only, allow all lower phase coefficients to vary:
\[
 f_{\mathbf c}(s)=\frac{s^{n+1}}{n+1}+\sum_{j=0}^n c_js^j,
 \quad h_{\mathbf c}=f_{\mathbf c}',\quad
 L_{\mathbf c}=u\partial_s+h_{\mathbf c},\quad
 \Pi(\mathbf c,u)_{j,b}
       =\int_{\Gamma_j}e^{f_{\mathbf c}/u}s^b\,ds.
 \tag{PDE3}
\]
The row index is $1\leq j\leq n$ and the column index $0\leq b<n$.
The original phase $f_t$ is the specified coefficient slice of this
family. No coefficient of that slice will be suppressed.

Ordinary monic division gives, for $0\leq j\leq n$ and $0\leq b<n$,
\[
 s^{b+j}=h_{\mathbf c}q_{b,j}+r_{b,j},\qquad
 \deg r_{b,j}<n,\quad\deg q_{b,j}\leq b+j-n.
 \tag{PDE4}
\]
A negative bound means the zero polynomial. In particular
$\deg q_{b,j}'\leq b-1<n$, so the exact differential reduction is
\[
 s^{b+j}=L_{\mathbf c}q_{b,j}
                    +(r_{b,j}-u q_{b,j}').
 \tag{PDE5}
\]
The leading term of $L_{\mathbf c}P$ has degree $n+\deg P$ and the same
leading coefficient as $P$. Repeated subtraction therefore proves that
the cokernel of $L_{\mathbf c}$ has the unique increasing frame of
degrees less than $n$, for every $\mathbf c$ and $u$. Formula PDE5 is
its literal remainder, not a specialization at $u=0$.

Let $B_j(\mathbf c,u)$ have these remainder columns, and let
$M_{\mathbf c}$ be multiplication by $s$ on
$\mathbb{C}[s]/(h_{\mathbf c})$. Ordinary remainders form
$M_{\mathbf c}^j$. The correction column $-u q_{b,j}'$ has row index
strictly less than $b$. It is consequently strictly triangular and has
zero trace. Thus, retaining its full nonzero matrix,
\[
 B_j=M_{\mathbf c}^{j}-u(\operatorname{coeff}q_{b,j}')_{b=0}^{n-1},
 \qquad \operatorname{Tr}B_j=\operatorname{Tr}M_{\mathbf c}^{j}.
 \tag{PDE6}
\]
The boundary integral of $e^{f_{\mathbf c}/u}q_{b,j}$ vanishes at
the two outer ends and cancels at the common origin. Differentiating
the original integrals and applying PDE5 gives
\[
 \partial_{c_j}\Pi=\frac1u\Pi B_j\quad(0\leq j\leq n).
 \tag{PDE7}
\]

Define the polynomial in the original coefficients
\[
 \mathcal S(\mathbf c)
  =\operatorname{Tr} f_{\mathbf c}(M_{\mathbf c})
  =\frac{\operatorname{Tr}M_{\mathbf c}^{n+1}}{n+1}
                +\sum_{j=0}^n c_j\operatorname{Tr}M_{\mathbf c}^j.
 \tag{PDE8}
\]
Its derivative with respect to $c_j$ is
$\operatorname{Tr}M_{\mathbf c}^j$. Here is a proof including repeated
roots. On the open set where $h_{\mathbf c}$ has distinct roots
$r_1,\ldots,r_n$, they admit local holomorphic enumerations and
$\mathcal S=\sum_i f_{\mathbf c}(r_i)$. Differentiation gives
$\partial_{c_j}\mathcal S=\sum_i r_i^j$, since
$f_{\mathbf c}'(r_i)=0$ makes every root-derivative term exactly zero.
Both sides of the identity are polynomials in $\mathbf c$; the open
set is dense because the derivative $s^n-1$ occurs in this family
and has distinct roots. Equality extends as a polynomial identity to
every coefficient value. Equivalently, the trace at a repeated root
is its full multiplicity times the scalar value, since powers of its
nilpotent part have zero diagonal. No multiplicity is removed.

Put $D=\det\Pi$. Multilinearity of the determinant in its columns
gives $\partial_{c_j}D=u^{-1}(\operatorname{Tr}B_j)D$ without any
prior assumption that $\Pi$ is invertible. PDE6 and PDE8 imply that
every derivative of $D\exp(-\mathcal S/u)$ is zero on the connected
coefficient space. Its value is determined at $\mathbf c=0$.

For completeness, write $\zeta=\exp(2\pi i/d)$ and
$\beta_b=(b+1)/d$. Direct integration on the original two rays gives
\[
 \Pi(0,u)_{j,b}
  =u^{(b+1)/d}d^{\beta_b-1}
          e^{i\pi\beta_b}\Gamma(\beta_b)
                       (\zeta^{j(b+1)}-1).
 \tag{PDE9}
\]
The power of $u$ uses the branch fixed in PDE2. To compute the remaining
determinant let $V_{r,b}=\zeta^{rb}$, $0\leq r,b<d$. Subtract row zero
from each other row and expand at its first column. The resulting
minor is exactly $(\zeta^{j(b+1)}-1)_{1\leq j\leq n,0\leq b<n}$,
with no row or column reversal. The Vandermonde formula gives
\[
 \det V=\prod_{0\leq p<q<d}(\zeta^q-\zeta^p),\qquad
 \zeta^q-\zeta^p
       =2i e^{\pi i(p+q)/d}\sin\frac{\pi(q-p)}d.
 \tag{PDE10}
\]
Every displayed sine is positive. The phase of this product is
$\pi d(d-1)/4+\pi(d-1)^2/2
=\pi(d-1)(3d-2)/4$.
Its first term comes from the $d(d-1)/2$ factors $i$, and its second
from $\sum_{p<q}(p+q)=d(d-1)^2/2$.
Also $V^*V=dI$ by the finite geometric sum, so $|\det V|=d^{d/2}$.
The retained gamma product is
$\prod_{b=0}^{n-1}\Gamma((b+1)/d)=(2\pi)^{n/2}/\sqrt d$,
the multiplication identity used in IPM21, also recorded in DLMF
5.5.7 at \url{https://dlmf.nist.gov/5.5.E7}.
The powers of $d$ from PDE9 sum to $-n/2$ and cancel this product's
$d^{-1/2}$ against $d^{d/2}$. Its additional phase is
$\pi\sum_b\beta_b=\pi n/2$. The total phase is therefore
$3\pi n(n+1)/4$. Consequently the exact determinant is
\[
 \boxed{\displaystyle
 \det\Pi(\mathbf c,u)
   =(2\pi)^{n/2}u^{n/2}
       \exp\left(\frac{3\pi i n(n+1)}4
                         +\frac{\mathcal S(\mathbf c)}u\right).}
 \tag{PDE11}
\]
In particular it never vanishes for $u\ne0$, including coefficients
with repeated critical points or repeated critical values. This proves
invertibility rather than assuming it to use a determinant equation.
For $n=1$ the formula is the downward Gaussian for $u>0$ on
$\arg u=0$, continued on the specified branch:
$-i\sqrt{2\pi u}\exp((c_0-c_1^2/2)/u)$.

\subsection{The original full quartet and the exact complex-parameter transport}

For the original full quartet $n=4m$, retain
$\rho=1/2+\delta+i\gamma$, $0<\delta<1/2$, $\gamma>2$, and
$h=\prod_{r\in\{\rho,\bar\rho,1-\rho,1-\bar\rho\}}(s-r)^m$.
Its reflection is $h(1-s)=h(s)$. With the original primitive constant
$A=\Phi_h(1/2)$, differentiation and evaluation at $s=1/2$ prove
\[
 \Phi_h(1-s)=2A-\Phi_h(s),\qquad
 f_t(1-s)=2A-t-f_t(s),\qquad
 \mathcal S(f_t)=n(A-t/2).
 \tag{PDE12}
\]
For the trace assertion, reflection pairs the roots of $h(s)-t$ with
their full multiplicities. Each pair has the displayed sum of phase
values. A root fixed by reflection has $s=1/2$ and phase value
$A-t/2$ itself. Summing over all $n$ roots, or taking the trace of the
corresponding algebra involution, proves the identity even at every
multiple-root parameter. Hence on this original quartet
\[
 \det\Pi_h(t,u)
  =(2\pi)^{n/2}u^{n/2}
       e^{3\pi i n(n+1)/4}\,e^{n(A-t/2)/u}.
 \tag{PDE13}
\]
Neither the original $A$ nor the $t$ term has been discarded.

